%==============================================================================
\documentclass[cic,tc]{iiufrgs}

\usetocstyle{standard}
\usepackage[utf8]{inputenc}   % acentuação
\usepackage{lmodern}          % allows arbitrary font size
\usepackage{graphicx}         % figuras
\usepackage{times}            % fonte Adobe Times
\usepackage{hyperref}
\usepackage[alf,abnt-emphasize=bf]{abntex2cite}	% cit abnt
\usepackage{mathtools}	      % math symbols
\usepackage{amsthm}           % theorems and definitionsi
\usepackage{listings}         % codeblocks

% chktex latex syntax correction tool configurations
% chktex-file 18 disable warnings about "
% chktex-file 19 disable warnings about ô
% chktex-file 6 disable italic correction warnings

%==============================================================================
\title{Classificação de imagens utilizando deep supervised learning}
\author{Ficagna}{Alan}
\advisor[Prof.~Dr.]{Engel}{Paulo Martins}
\date{}{2015}
\keyword{IA}
\keyword{inteligência artificial}
\keyword{deeplearning}

%==============================================================================

\graphicspath{{fig/}}

\DeclareGraphicsExtensions{.pdf,.png,.jpg}

\begin{document}

\maketitle

\clearpage
\begin{flushright}
  \mbox{}\vfill
  {\sffamily\itshape{}
    ``Virtue is more to be feared than vice,\\
    because its excesses are not subject to the regulation of consciousness''\\}
  --- \textsc{Adam Smith}
\end{flushright}

%==============================================================================
\begin{abstract}

Este documento é uma revisão do estado da arte em relação as técnicas de
Deep Learning utilizadas no reconhecimento de caracteres escritos à mão, tendo o
banco de dados MNIST como meio de comparação. Além disso são apresentadas
algumas variações nos meta parâmetros do algoritmo tais como <TODO> para
avaliar o impacto no tempo de execução e na qualidade do resultado.
\end{abstract}

\begin{englishabstract}{Recognizing Handwritten Characters Using Deep Learning Algorithms}{IA, deeplearning} % tittle, keywords

This document is a review of the state of the art in regards to Deep Learning
techniques used in handwritten characters recognition having the MNIST
database for comparison. Also, the impact in the change of parameters such as
<TODO> was measured to provide and indication of its effects on run speed and
quality of result.
\end{englishabstract}

%==============================================================================
\listoffigures
\listoftables
\begin{listofabbrv}{ILSVRC} % largest
 \item[MNIST] Mixed National Institute of Standards and Technology
 \item[CNN] Convolutional Neural Network
 \item[GPU] Graphical Processing Unit
 \item[ReLU] Rectified Linear Unit
 \item[CONV] Convolutional
 \item[FC] Fully Connected
 \item[ILSVRC] Large Scale Visual Recognition Challenge
\end{listofabbrv}
\tableofcontents

%==============================================================================

\chapter{Introdução}

Nos últimos anos houve um grande avanço na Computação Visual graças ao
ressurgimento de um campo da Inteligência Artificial conhecido como Redes
Neurais. Este progresso foi responsável pelo despertar do interesse e
investimentos de grandes empresas como Google, Facebook e outros que apostam no
desenvolvimento de aplicações que há poucas décadas atrás eram apenas tema de
ficção científica tais como scanners faciais de alta precisão, reconhecimento
de objetos em cenas, carros que não precisam de motorista e vários tipos de
classificadores em geral. Essas novas técnicas ficaram conhecidas como Deep
Learning e prometem fazer uso da enorme quantidade de dados e informações
disponíveis na internet e nas bases de dados das empresas.

Este progresso foi causado por fatores como o incremento da capacidade de
processamento dos computadores e a ubiquidade de tecnologias como GPUs que
facilitaram o processamento vetorial e paralelo de imagens, sons e outros
domínios de alta dimensionalidade que até então possuíam um custo computacional
proibitivo.

\section{Objetivos}

<REVISAR>
O objetivo deste trabalho consiste em efetuar uma revisão do
histórico do desenvolvimento das técnicas de Deep Learning e um resumo do seu
estado da arte com a finalidade de adquirir proficiência no uso das mesmas.
Além disso, serão apresentados resultados práticos obtidos com a execução de
uma modelo que utiliza Deep Learning chamado de Rede Neural Convolucional
(CNN) no reconhecimento de caracteres escritos à mão utilizando o algoritmo da
rede LeNet de LeCun et al sobre a base MNIST.\@



% chktex-file 18 disable warnings about "
% chktex-file 19 disable warnings about ô
% chktex-file 6 disable italic correction warnings

\chapter{Histórico}

\section{Inteligência Artificial}

Definir o que é inteligência não é uma tarefa fácil, e inteligência artificial
é mais difícil ainda. Historicamente, existiram duas abordagens para a
definição de inteligência: O comparativo com humanos e o comparativo com seres
racionais.

\subsection{Inteligência Humana}

No primeiro modelo, o seres da espécie humana são tidos como parâmetros e
portanto, qualquer avaliação de inteligência tem que se refletir em
comportamentos ou pensamentos similares aos humanos. O Teste de Turing proposto
pelo matemático Alan Turing (1950) foi uma maneira operacional de estabelecer
inteligência pelo comparativo comportamental. Em sua versão mais simplificada,
um computador poderia ser considerado inteligente se, ao ser interrogado por um
humano com algumas questões, este não pudesse determinar se as respostas eram
oriundas de um computador ou de uma pessoa. A Ciência Cognitiva, por outro
lado, busca compreender o mecanismo pelos quais os humanos são capazes de
pensar e à partir disso estabelecer comparações com algum algoritmo que
eventualmente pudesse replicar estas capacidades.

\subsection{Inteligência Racional}

Além do paradigma humano, existe um outro que é a comparação com seres
racionais. Um ser é dito racional se ele age da ``maneira correta'', dado o que
ele conhece. Isso nos leva a seguinte definição: um ser é inteligente se ele
age ou pensa como um ser racional.

Aristóteles foi o primeiro a sistematizar uma maneira de ``pensar
corretamente''.  Ele foi o criador da disciplina que hoje conhecemos por Lógica
que estabelecia estruturas argumentativas tais como: ``Sócrates é um homem;
todos os homens são mortais; portanto, Sócrates é mortal'' as quais sempre
produzem conclusões corretas. Porém, estes silogismas nem sempre são capazes de
traduzir conhecimento informal, principalmente em face à incertezas e por isso
tendem a não ser reconhecidos como tudo que existe em termos de inteligência.

Além disso existem diferenças entre resolver um problema em princípio e
resolvê-lo na prática. Embora a abordagem pelas ``leis do pensamento'' do
Aristóteles foquem em realizar inferências corretas, esta não é capaz de
exaurir o que significa agir como um ser racional pois existem situações em que
não há uma ação comprovadamente correta, mas em que uma ação deve ser tomada
mesmo assim. Além disso, existem comportamentos racionais que não podem ser
deduzidos por mera inferência como, por exemplo, o reflexo de remover
rapidamente a mão do fogo sem uma cuidadosa pré deliberação sobre o assunto.

O propósito deste trabalho não é exaurir todos os ângulos pelos quais esta
questão já foi abordada, tais como psicologia, lógica e filosofia mas sim
apresentar apenas uma introdução que permita entender o surgimento do o campo da
Inteligência Artificial dentro da ciência da computação.

\section{Histórico da Inteligência Artificial}

\subsection{Primeiros anos}

O primeiro trabalho a ser considerado efetivamente ``Inteligência Artificial''
foi feito em 1943 por Warren McCulloch e Walter Pits, no qual foi proposto o
modelo de funcionamento de um neurônio artificial.  Este neurônio podia assumir
os estados ``ligado'' e ``desligado'', determinado em resposta aos estímulos
proporcionados pelos outros neurônios vizinhos. Foi demonstrado também que uma
rede de neurônios artificial conseguia computar qualquer função computável e
sugerido que a mesma poderia aprender.

Em 1949 Donald Hebb demonstrou uma simples regra pela qual a força ou os pesos
numéricos de conexão entre os neurônios poderia ser modificada, a qual ficou
conhecida como aprendizado hebbiano. Logo em seguida, em 1950, foi construído o
primeiro computador de rede neural chamado de SNARC que fora feito com 3000
tubos de vácuo e conseguia simular uma rede de 40 neurônios.

Embora o neurônio artificial seja considerado o primeiro trabalho, o termo
Inteligência Artificial só viria a ser cunhado mais de 10 anos depois em
Darthmouth nos Estados Unidos em um workshop organizado por John McCarthy que
duraria dois meses e cujo propósito era estudar teoria dos autômatos, redes
neurais e inteligência. Embora não tenha produzido nenhum progresso, o workshop
serviu para introduzir as principais figuras da área entre si.

Nos seus primeiros anos a IA obteve bastante sucesso. Um a um, os items da
lista ``uma máquina nunca vai poder fazer X'' compilada por Turing iam sendo
desbancados. Newel and Simon criaram o primeiro GPS (General Problem Solver),
um programa desenhado para imitar a maneira como humanos pensam ao invés de
simplesmente seguir regras lógicas. Na classe limitada de desafios que
conseguia resolver, a ordem dos subobjetivos que o programa tentava seguir era
similar à humana. Na IBM, Nathaniel Rochester criaria o primeiro Provador de
Teoremas Geométricos, o qual conseguia provar teoremas que vários estudantes
consideravam complexos.  Arthur Samuel refutou a ideia de que programas somente
conseguiriam fazer o que lhes eram dito criando um programa de computador que
conseguia jogar um jogo melhor que o seu criador.

Foi nessa época também que McCarthy criaria sua linguagem de programação
\emph{Lisp} que viraria a linguagem dominante em IA e com a qual ele faria o
primeiro algoritmo que aceitava novos axiomas, sendo então o primeiro autômato
que conseguia adquirir novas capacidades sem necessitar ser reprogramado.
Outras contribuições significativas foram o programa SAINT (1963) que conseguia
resolver problemas de Cálculo típicos do primeiro ano de curso, o programa
ANALOGY (1968) respondia a perguntas de analogia geométrica como as de testes
de QI, o programa STUDENT (1967) resolvia problemas algébricos e o programa
SHRDLU (1972) era capaz de executar instruções do tipo ``Encontre um bloco
maior do que você está segurando agora e o coloque na caixa''. Foi nessa época
também que foram feitas melhorias no processo de aprendizado Hebbiano que
ficaram conhecidas como redes \emph{adalines}.

\subsection{Maturação e Estado da Arte}

O sucesso inicial foi seguido de um período de relativa estagnação. O fato de
que estes algoritmos só conseguiam fazer manipulação simbólica e não conheciam
nada sobre o domínio que trabalhavam provou ser um grande empecilho. Além disso
a maioria dos primeiros programas de IA resolviam o problema tentando várias
combinações até que o resultado fosse atingido, o que não era factível para
instâncias de problemas marginalmente maiores. Foi nesta época que Minsky e
Papert provaram que um neurônio artificial agora conhecido como
\emph{perceptron} era restrito no tipo de funções que conseguia representar e o
financiamento de pesquisas na area se tornou escasso. Já na década de 1970,
houve o surgimento dos sistemas baseados em conhecimento prévio como o programa
DENDRAL que viu uma melhora significativa de desempenho na tarefa de inferir a
estrutura molecular à partir da massa dos seus componentes. O programa que
inicialmente tentava todas as alternativas possíveis passou a utilizar de
conhecimento de especialistas em química para identificar subestruturas
conhecidas o que reduzia significadamente o número de tentativas.

O primeiro caso de sucesso comercial reportado destes sistemas especialistas
era um programa que ajudava na compra de novos computadores e que teria
economizado à empresa US\$40 milhões por ano, o que explica o fato destes
sistemas continuarem sendo usados até hoje.

Paralelamente, houve o ressurgimento das redes neurais no final da década de
1980, proporcionado pela reinvenção do algoritmo de \emph{back-propagation},
uma maneira eficiente de atualizar os pesos das conexões nos neurônios, e
também uma melhoria nas práticas de pesquisa na área que passou a exigir um
maior rigor científico acompanhado do uso de bancos de dados padrões que
permitiam o teste e comparação dos resultados entre diferentes algoritmos. Em
uma outra linha recente, a ênfase no algoritmo da lugar ao uso de grandes
quantidades de dados inspirado no trabalho de Yarowsky (1995) que demostrou ser
possível reconhecer se o uso da palavra \emph{plant} significava uma flor ou uma
fábrica com precisão acima de 96\% sem utilizar de rótulos fornecidos por
humanos, apenas com as definições do dicionário desta palavra. Trabalhos como o
dele sugerem um ``gargalo de conhecimento'' na IA, ou seja, o conhecimento que
o um sistema precisa pra resolver um problema pode vir de grandes massas de
dados ao invés de intervenções de especialistas na área. Hoje em dia, a IA está
em várias aplicações de campos diferentes, como por exemplo:

\begin{itemize}
\itemsep1pt\parskip0pt\parsep0pt
\item
  Na área de robótica veicular com carros auto guiados.
\item
  Na área de reconhecimento de fala.
\item
  Planejamento automático de tarefas.
\item
  Jogos, com destaque para o computador DEEP BLUE da IBM que derrotou o
  campeão mundial de xadrez, Garry Kasparov.
\item
  Controle de SPAN em mensagens de e-mail.
\item
  Militar e Logística.
\end{itemize}



% chktex-file 18 disable warnings about "
% chktex-file 19 disable warnings about ô
% chktex-file 6 disable italic correction warnings


\chapter{Revisão sobre Redes Neurais}

As primeiras definições de uma rede neural artificial foram feitas em 1943 por
McCulloch e Pits inspiradas na hipótese de que as atividades mentais do cérebro
humano consistem em reações eletroquímicas em redes de células cerebrais
chamadas de neurônios. O neurônio artificial consiste em uma unidade
computacional abstrata capaz de receber estímulos de entrada e produzir uma ou
mais saídas e uma rede neural artificial --- de agora em diante referenciada
apenas por \emph{rede neural} --- consiste em um grafo direcionado cujos nodos
são neurônios artificiais.

\begin{figure}
  \caption{Abstração de um neurônio artificial}
  \begin{center}
    \includegraphics[scale=0.5]{placeholder}
  \end{center}
\end{figure}

\section{O neurônio artificial}

Cada neurônio recebe $m$ entradas ponderadas por um peso $w$,
correspondente a entrada do seu i-ésimo vizinho. Esta configuração foi
inspirada nos neurônios naturais cujas interconexões são reguladas por um
potencial de ativação responsável por diferenciar a intensidade entre as
mesmas. A saída é produzida aplicando-se a função limite $\sigma$ na soma
de todos as entradas ponderadas do neurônio. Esta função simula a \emph{lógica
  de limite} existente nos neurônios naturais, os quais somente irão propagar o
seu sinal caso haja um acúmulo suficiente de potencial elétrico no mesmo. Em
sua forma matemática, a saída $y$ do neurônio artificial pode ser expressa da
seguinte forma:

$$y=\sigma(\sum_{i=1}^{m}x_i w_i + b)$$

$\sigma$ representa uma classe de funções que podem ser utilizadas dependendo
do objetivo da construção da rede neural. Para problemas de classificação
\emph{linearmente separáveis}, pode ser usada a função limite exemplificada na
figura~\ref{fig:lim-fun}, enquanto para problemas não linearmente separáveis é
necessário utilizar algum tipo de função \emph{sigmoide}. A função logística é
um exemplo de função sigmoide frequentemente utilizada por ser diferenciável em
todo o seu domínio, o que permite produzir garantias de convergência nos
algoritmos de aprendizagem.

O valor $b$ corresponde ao \emph{bias} aplicado a função e pode ser omitido
simplesmente inserindo-se artificialmente uma nova entrada $x_0 = 1$.

\begin{figure}\label{fig:lim-fun}
  \caption{Exemplo de funções de limite $\sigma$ usadas em neurônios artificiais}
  \begin{center}
    \includegraphics[scale=0.5]{placeholder}
  \end{center}
\end{figure}

\subsection{Bias}

Todo neurônio artificial é dito ter pelo menos uma entrada extra $x_0 = 1$ e
consequentemente, o peso extra associado $w_0$ chamado de bias $b$. O bias
permite que a função aproximada pelo neurônio não necessite obrigatoriamente
passar pela origem $(0,0)$, o que reduz o erro médio da aproximação, caso os
dados de treinamento não tenham sofrido nenhum tipo de ajuste numérico, como por
exemplo, a remoção do valor médio.

\begin{figure}
  \caption{Exemplo de função aproximada com e sem o vetor de bias}
  \begin{center}
    \includegraphics[scale=0.5]{placeholder}
  \end{center}
\end{figure}

\subsection{Redes neurais}

Em uma rede neural, os neurônios são organizados em \emph{camadas}. Alguns
autores não consideram a camada de entrada como uma camada distinta, mas esta
abstração permite simplificar a notação e torna o processo de conectar redes
neurais entre si mais intuitivo. Outra camada especial é a de saída, pois é
ela que determina se a rede exercerá a função de um \emph{classificador} ou de
um \emph{regressor}. Todas as demais camadas intermediárias são chamadas de
\emph{camadas ocultas}; o nome não possui nenhuma conotação especial e serve
apenas para diferenciá-las das camadas de entrada e saída. Em uma rede
\emph{feed forward} as saídas de cada camada somente são computadas na direção
e sentido da camada de saída, isto é, não há uma retroalimentação das conexões
de saída para as de entrada. Quando isto acontece a rede é chamada de
\emph{recorrente} porém, o estudo deste tipo de rede neural não esta nos
objetivos deste trabalho.

Sendo assim, uma rede neural necessita ter pelo menos duas camadas --- a de
entrada e a de saída. A adição de camadas ou neurônios não possui um efeito
óbvio e generalizável para qualquer rede, e portanto é alvo de constante
experimentação e observação empírica sujeito a comportamento variável
dependendo do problema.

\begin{figure}\label{fig:003-nn}
  \caption{Exemplo de uma rede neural de 3 camadas}
  \begin{center}
    \includegraphics[scale=0.5]{placeholder}
  \end{center}
\end{figure}

Cada camada $S$ é formada por $R$ neurônios, cujos pesos das conexões podem ser
representados através de uma matriz $\mathbf{W}$. A notação em forma vetorial e
matricial dos pesos das camadas permite que operações mais eficientes sejam
utilizadas para realizar a computação da saída da rede neural e do ajuste dos
mesmos.

$$ W_{S,R} =
\begin{bmatrix}
  w_{1,1} & w_{1,2} & \cdots & w_{1,R} \\
  w_{2,1} & w_{2,2} & \cdots & w_{2,R} \\
  \vdots  & \vdots  & \ddots & \vdots  \\
  w_{S,1} & w_{S,2} & \cdots & w_{S,R}
\end{bmatrix}
$$

A saída ou \emph{ativação} de uma camada $\mathbf{\vec a^{\prime}}$ pode ser expressa em função da
ativação da camada anterior $\mathbf{\vec a}$ da seguinte maneira:

$$\mathbf{\vec a^{\prime}} = \sigma(\mathbf{\vec W\vec a+\vec b})$$

\subsection{Aprendizagem}

A aprendizagem em uma rede neural se da pelo ajuste dos pesos $w$ das conexões
dos neurônios. Ao longo dos anos foram propostas várias maneiras de atualizar
os pesos de uma rede neural automaticamente através de uma \emph{regra de
  aprendizagem}. Estas regras foram categorizadas em três grandes grupos:

\begin{description}

  \itemsep1pt\parskip0pt\parsep0pt

  \item[Aprendizagem Supervisionada] As entradas são acompanhadas por um
    conjunto de valores corretamente associados a mesma.

  \item[Aprendizagem Por Reforço] Semelhante ao aprendizado supervisionado,
    porém ao invés da resposta correta, é fornecido algo como uma nota ou uma
    indicação da performance da predição.

  \item[Aprendizagem Não Supervisionada] Nenhuma informação adicional a
    respeito das entradas é fornecida ao algoritmo.

\end{description}

\begin{figure}\label{fig:limit-functions-exapmple}
  \caption{Exemplo de funções de ativação}
  \begin{center}
    \includegraphics[scale=0.5]{placeholder}
  \end{center}
  \legend{Na esquerda uma simples função de limite; à direita a função $\frac{1}{1+e^{-x}}$}
\end{figure}

\subsubsection{Backpropagation}

Obviamente que a aplicabilidade de redes neurais seria extremamente limitada se
não houvesse uma maneira eficiente de aprender os pesos da rede de forma
automática.  Atualmente, a melhor maneira conhecida de fazer o mesmo é através
do algoritmo de \emph{backpropagation}. Este algoritmo embora já sabido desde a
década de 60, ganhou visibilidade na área de IA na década de 80, quando
aplicado em redes neurais.\citep{russell1995modern}

O processo de aprendizagem de uma rede neural segue a noção intuitiva que temos
do assunto: Um individuo faz uma predição acerca de uma proposição ou pergunta,
avalia o \emph{feedback} recebido do ambiente e corrige a sua previsão baseado
na distância entre o que foi previsto e o que de fato aconteceu.

A determinação da direção no espaço de pesos da correção é dado pelo
\emph{gradiente} função de erro $\mathbf{L}$, ou como é mais comumente
conhecida, \emph{loss function}. A eficiência do algoritmo de backpropagation
se baseia no fato de que a computação de cada camada depende somente da
computação das camadas anteriores.

Embora o algoritmo use a derivada da \emph{loss function} pré calculada de
maneira analítica, é comum ver nas implementações o uso da computação do
gradiente de maneira iterativa como forma de validação da implementação.

A computação completa da rede neural então, se da em dois passos: O primeiro
chamado de \emph{forward pass}, a ativação de cada camada é computada até que
se chegue numa resposta na camada de saída, a qual é usada para calcular o erro
em cada neurônio e o ajuste dos pesos $\Delta W_{i,j}$ do mesmo de acordo com a
seguinte equação:

$$ \mathbf{\Delta W_{i,j} = -\alpha\frac{\partial L}{\partial W_{i,j}}} $$

Onde $\alpha$ representa a taxa de aprendizagem.


\section{Why deep architetures}

% notes from lecun's https://www.youtube.com/watch?v=oOB4evKlEmQ

% Why deep learning instead of shallow learning since any deep architeture can
% be simulated by a shallow one?  According to yann le'cun it's because that is
% a trade-off between 'space' and computing time. It's true that shallow
% architetures can be simulated by a deeper one, but they run much much slower.

% See: Scaling learning algorithms towards AI

% Deep Supervised learning is hard beacuse the loss functions that are used in
% traing are non convex

% Backpropagation does not work well with networks that are tall and skinny
% lots of layers with few hidden units

% How to deal with limitations of linear classifiers You map the input to a
% higher dimension input where they are more easily separable

% How to build this: Use knowlege you already have to craft a feature set. We
% use a standard set of basis functions RBF Simplest approach use a kernel
% method


\input{tex/demonstration-lenet}

\section{Conclusão}

<TODO>


%==============================================================================

\bibliographystyle{abntex2-alf}
\bibliography{biblio}

\end{document}
